% Generated from the official PLOS template version 3.8, April 2026.
\documentclass[10pt,letterpaper]{article}
\usepackage[top=0.85in,left=2.75in,footskip=0.75in]{geometry}

% amsmath and amssymb packages, useful for mathematical formulas and symbols
\usepackage{amsmath,amssymb}

% Use adjustwidth environment to exceed column width (see example table in text)
\usepackage{changepage}

% textcomp package and marvosym package for additional characters
\usepackage{textcomp,marvosym}

% cite package, to clean up citations in the main text. Do not remove.
\usepackage{cite}

% Use nameref to cite supporting information files (see Supporting Information section for more info)
\usepackage{nameref,hyperref}

% line numbers
\usepackage[right]{lineno}

% ligatures disabled
\usepackage[nopatch=eqnum]{microtype}
\DisableLigatures[f]{encoding = *, family = * }

% color can be used to apply background shading to table cells only
\usepackage[table]{xcolor}

% array package and thick rules for tables
\usepackage{array}

% create "+" rule type for thick vertical lines
\newcolumntype{+}{!{\vrule width 2pt}}

% create \thickcline for thick horizontal lines of variable length
\newlength\savedwidth
\newcommand\thickcline[1]{%
  \noalign{\global\savedwidth\arrayrulewidth\global\arrayrulewidth 2pt}%
  \cline{#1}%
  \noalign{\vskip\arrayrulewidth}%
  \noalign{\global\arrayrulewidth\savedwidth}%
}

% \thickhline command for thick horizontal lines that span the table
\newcommand\thickhline{\noalign{\global\savedwidth\arrayrulewidth\global\arrayrulewidth 2pt}%
\hline
\noalign{\global\arrayrulewidth\savedwidth}}


% Remove comment for double spacing
\usepackage{setspace}
\doublespacing

% Text layout
\raggedright
\setlength{\parindent}{0.5cm}
\textwidth 5.25in
\textheight 8.75in

% Bold the 'Figure #' in the caption and separate it from the title/caption with a period
% Captions will be left justified
\usepackage[aboveskip=1pt,labelfont=bf,labelsep=period,justification=raggedright,singlelinecheck=off]{caption}
\renewcommand{\figurename}{Fig}

% Please use the included `plos2025.bst` as your BibTeX style
\bibliographystyle{vendor/plos2025}

% Remove brackets from numbering in List of References
\makeatletter
\renewcommand{\@biblabel}[1]{\quad#1.}
\makeatother



% Header and Footer with logo
\usepackage{lastpage,fancyhdr,graphicx}
\usepackage{epstopdf}
%\pagestyle{myheadings}
\pagestyle{fancy}
\fancyhf{}
%\setlength{\headheight}{27.023pt}
%\lhead{\includegraphics[width=2.0in]{PLOS-submission.eps}}
\rfoot{\thepage/\pageref{LastPage}}
\renewcommand{\headrulewidth}{0pt}
\renewcommand{\footrule}{\hrule height 2pt \vspace{2mm}}
\fancyheadoffset[L]{2.25in}
\fancyfootoffset[L]{2.25in}
\lfoot{16 November 2026}

%% Include all user-defined macros below

\newcommand{\lorem}{\textbf{LOREM}}
\newcommand{\ipsum}{\textbf{IPSUM}}

%% END MACROS SECTION



\usepackage[T1]{fontenc}
\usepackage[utf8]{inputenc}
\usepackage{longtable}
\hypersetup{hidelinks}
\setlength{\emergencystretch}{2em}


\begin{document}

\vspace*{0.2in}

\begin{flushleft}

{\Large\textbf{Equivalent names yield inconsistent factual answers in two language models}}\par\bigskip

Ihor Kendiukhov\textsuperscript{1*}\par\medskip

\textbf{1} Institute of Medical Genetics and Applied Genomics, University of Tübingen, Tübingen, Germany\par\medskip

* ihor.kendiukhov@student.uni-tuebingen.de\par

ORCID: \url{https://orcid.org/0000-0001-5342-1499}\par

Short title: Equivalent names and inconsistent factual retrieval

\end{flushleft}

\clearpage

\newgeometry{top=0.85in,left=1in,right=1in,footskip=0.75in}

\fancyheadoffset[L]{0in}\fancyfootoffset[L]{0in}

\section*{Abstract}

Two names explicitly assigned to the same object should retrieve the same facts. We examined this requirement in two instruction-tuned language models using fictional registries whose facts and name bindings were supplied entirely in context. Thirty-two counterbalanced registries generated 3,072 questions per model, including 2,304 factual questions. Qwen3-4B-Instruct-2507 answered 1,657 factual questions correctly and produced 647 errors, all of which selected another listed object's value for the requested field. Phi-3.5-mini-instruct answered 1,980 correctly, selected another object's value in 260 cases, and produced 64 other factual errors. Equivalent-name pairs with one correct answer and one completed wrong-object answer occurred in 23 and 14 registries, respectively. In Qwen, transferring a correct-name attention allocation across fields initially corrected five of eight first-token predictions. The fixed intervention subsequently repaired only one of eight complete-answer errors in six other development registries while preserving ten initially correct answers. A nine-condition diagnosis on those same six registries tested record allocation, within-record attention, record contributions, nonrecord contributions, whole-component transfer, and controls. None of the six new targeted variants repaired an error. Whole-component transfer nevertheless improved the mean correct-minus-original-wrong token score by 4.214, demonstrating that an operative, directionally favorable intervention need not cross the answer decision boundary. Numerical replay and independent literal-prompt grading agreed with all reported counts. These results establish reference inconsistency in the tested settings and constrain one proposed repair interface. They do not identify a new binding mechanism or establish that entity representations are absent. All causal follow-ups remain development analyses; reserved confirmation registries were not evaluated.

\clearpage

\newgeometry{top=0.85in,left=1in,right=1in,footskip=0.75in}

\fancyheadoffset[L]{0in}\fancyfootoffset[L]{0in}

\linenumbers

\section*{Introduction}

A language model using an explicitly defined registry must connect a requested name to the right object before returning that object's properties. If two names denote the same object, changing between them should not change the answer to an otherwise identical factual question. This requirement concerns consistency under an equivalence relation supplied by the prompt. It can fail even when the required facts are present, the answer vocabulary is small, and neither name has a unique real-world referent.


Related failures and mechanisms are already documented. Yang showed that familiar programming-library aliases can override explicit bindings and that correct library identification need not ensure correct attribute use \cite{yang2026import}. Gur-Arieh and colleagues distinguished lexical, positional, and direct-pointer retrieval mechanisms in contextual binding tasks \cite{gurarieh2026mixing}. Query-conditioned rebinding circuits and structured entity-by-relation representations have also been reported \cite{ohdemberg2026retrieval,dai2026cell}. These findings motivate distinguishing a detectable reading pattern from an intervention that transfers reliably to other cases. We ask whether an apparently informative record-allocation pattern provides such a repair for explicit reference errors.


A repair claim needs more than a change in an internal representation or an improvement in an output score. Interventions should produce the intended answer while preserving other answers, as emphasized by attribute-intervention benchmarks such as RAVEL \cite{huang2024ravel}. The wider debate over attention and explanation likewise distinguishes a readable weighting pattern from an adequate account of a model's computation \cite{jainwallace2019attention,wiegreffepinter2019attention}. Activation-patching results depend on the comparison, intervention location, and evaluation metric \cite{zhangnanda2024patching}. A causal abstraction requires that the proposed internal-to-conceptual mapping support the relevant interventions, rather than merely correlate with behavior \cite{geiger2025causal}.


We studied two instruction-tuned model families using artificial registries with four objects, three factual fields, and explicit aliases. Two aliases referred to one object. We balanced object identifiers, alias roles, list positions, and wording, and crossed consistent exchanges of the equivalent names with highlighted-alias conditions. We then tested a narrowly selected component in one model. The intervention transferred a correct-name record-allocation profile while retaining the recipient's own record contents, making it more restrictive than copying a donor answer.


The behavioral finding replicated across the two models. The proposed causal repair did not. A promising discovery result weakened on unused development registries, and a subsequent fixed-component decomposition found no reliable repair despite substantial changes in prediction scores. Our contribution is a controlled empirical account of this reference-consistency failure and of the limits of this particular repair hypothesis. We do not propose a new general binding mechanism. The study separates behavioral replication, exploratory component selection, development validation, and reopened diagnosis; none is presented as final independent confirmation.


\section*{Materials and methods}

\subsection*{Study design and units of analysis}

The evaluated dataset comprised 32 newly generated development registries, hereafter called worlds. Each world specified four objects with permanent identifiers A through D, a display nameplate, and distinct values for color, shape, and home. Four aliases were assigned to objects: two aliases referred to a common target object, two referred to different objects, and one object had no alias. The complete ordered assignment of values to objects was unique for each field across the generated study, although individual words recurred. The task was contextual lookup, not retrieval of memorized facts about actual countries, people, or companies.


Within each world, two assignments exchanged the equivalent names between their binding-line positions. Because both names denoted the same object, the word-to-object mapping and all factual answers remained unchanged. Each assignment was crossed with four highlighting conditions: no highlighted alias, the first equivalent alias, the second equivalent alias, or an unrelated alias. Each of the resulting eight contexts had nine factual questions, two object-identifier questions, and one global status-word question. Thus each world contributed 96 questions, including 72 facts, 16 identifiers, and eight status questions; each model received 3,072 questions, including 2,304 facts.


Two wording formats were balanced between worlds. Target record position and the first equivalent alias's binding-line position covered all 16 combinations equally often. Object identifiers, lexical names in semantic alias roles, and the two equivalent-name line orders were also balanced as specified by the generator. Remaining combinations were randomized rather than claimed to form a complete factorial design. The alias bank was apple, bat, anchor, date, cat, basket, fig, and dog. These common words were selected using the two pinned tokenizers before model outputs were observed. This design tests context-defined alias equivalence; it is not a general evaluation of natural-language word-sense disambiguation.


The independent dataset unit was a world. Name exchanges, highlighting conditions, fields, and repeated model queries were within-world observations. Both models evaluated the same 32 worlds, not separate samples totaling 64 independent worlds. A further 64 worlds were generated and tokenized as reserved confirmation material but received no model forward passes. Earlier project experiments using different tasks or smaller pilots are outside the present study.


\subsection*{Models and deterministic inference}

We evaluated Qwen3-4B-Instruct-2507 and Phi-3.5-mini-instruct \cite{qwenmodel,phimodel}. Model revisions, token assets, prompts, and execution sources were hash-pinned. Both models ran in evaluation mode with frozen parameters, 32-bit floating-point arithmetic, eager attention, batch size one, and TensorFloat-32 disabled on an NVIDIA A40. The runtime used Python 3.12.15, PyTorch 2.6.0 with CUDA 12.4, and Transformers 5.14.1. Exact revisions and additional runtime settings are given in S1 Appendix and the accompanying records. These settings prioritize reproducibility and are not intended to characterize optimized deployment throughput.


Each prompt instructed the model to return only the requested value. The requested alias preceded an explicit boundary marker, and the requested field followed it. Greedy generation stopped only on a configured end-of-sequence token or after eight generated tokens. Decoding removed only a final end-of-sequence token and retained other special tokens. A completed answer was correct only if its entire decoded text matched the literal answer after case, whitespace, and final-punctuation normalization. Explanations, prefixed answers, and truncations were not silently shortened into correct responses.


\subsection*{Behavioral outcomes}

A factual wrong-object answer was a completed response matching another object's listed value for the requested field. Distinct field values made this classification unambiguous. Other incorrect responses, including malformed or truncated answers, were retained separately. For each world, name assignment, highlighting condition, and field, we paired the two equivalent-name questions. An eligible wrong-object disagreement required one correct response and one completed response selecting another listed object's value. Each world supplied 24 such pairs, giving 768 pairs per model. Broader string disagreement and any-one-correct disagreement were secondary outcomes and were not conflated with the narrower eligible category.


The model-specific feasibility criterion required wrong-object facts in at least 12 worlds and at least 24 eligible equivalent-name disagreements distributed across at least eight worlds. This criterion selected a viable subject for mechanism investigation; it was not a significance test or a criterion for novelty. Identifier probes were scored under the same strict whole-answer rule. We also inspected identifier formatting descriptively, without replacing the primary score with a more permissive one.


\subsection*{Staged causal investigation}

Only Qwen received internal interventions. Eligible contexts required a donor alias that answered all three facts and the identifier correctly, together with an equivalent recipient alias producing at least one completed wrong-object fact. Deterministic hash ordering selected one context per eligible world. There were 14 eligible worlds: eight formed the discovery cohort and six were initially unused for intervention selection. All were drawn from the behaviorally evaluated development set. Subsequent reuse of the six-world cohort was explicitly exploratory, not a second validation sample.


An initial bounded screen replaced attention outputs at the alias-boundary position before any field was supplied. It screened nonfinal layers and then heads in two selected layers. Effects were weak and did not justify a completed-answer repair claim. We next examined attention at the actual answer query. An attention head is a component that weights earlier token positions and combines their value vectors into an output. Natural captures contrasted the correct and incorrect equivalent-name computations, using both color and a hash-selected erroneous field in each discovery case.


Record regions were obtained from the printed registry lines and tokenizer offsets, without using the correct object's identity to construct a region. For each head, attention was summarized as the total mass assigned to all records, a four-record profile normalized within that total, and a distribution over tokens within each record. A record-profile intervention copied all four allocation proportions from a source computation. It retained the recipient's total record attention, relative token weights within each record, all nonrecord weights, and original value vectors. Correct- and wrong-object labels were used for exploratory selection and scoring, but were not inputs to the numerical transfer operation. The source color profile was reused for other requested fields, so it did not transfer the requested field's answer vector.


Natural attention contrasts across all 36 layers and 32 heads, or 1,152 heads, were ranked using the eight labeled discovery cases before the top two supported positive candidates were tested. A fixed ordering based on the advantage over a cross-world source-profile control selected decoder layer 28, head 0, using zero-based indices, for subsequent testing. The operation was then tested unchanged on the six remaining eligible development worlds across all three recipient fields. They contained eight initially wrong and ten initially correct facts. Controls included self transfer, a frozen discovery-average profile, a same-alias color profile, reverse transfer, and cyclic record reassignment. Repair and predicted wrong-object steering were scored separately; counterfactual target hits were not counted as factual accuracy.


The fixed development criterion required at least four of eight errors to be repaired in at least three worlds, no more than one of ten initially correct facts to be damaged, and at least two more repairs than the frozen-average control. Reverse and cyclic steering had separate targets of at least nine of 18 specified outcomes across at least three worlds; cyclic steering additionally required at least two predicted object identifiers. These targets were fixed before those intervention outcomes. Their failure stopped this candidate's validation route.


\subsection*{Fixed component diagnosis}

A separately frozen diagnosis reopened the same six development worlds after the repair failure. It retained the component, cases, decoding rule, and intervention strength. For each field it captured the actual head output, attention weights, and recipient value vectors. Record value tensors were byte-identical across the six natural alias-by-field contexts within each case, as expected for a shared causal prefix. This invariance is an architectural control, not evidence of an abstract entity code.


Nine conditions were evaluated: an unchanged self control; the original color-derived record profile; the donor's same-field record profile; its within-record token pattern; its profile and within-record pattern jointly at the recipient's record total; its entire record contribution; its nonrecord contribution; its complete head output; and a random direction matched to the complete-head change's norm. The record contribution was the attention-weighted sum of value vectors at record positions. The nonrecord contribution was defined as the actual head output minus that record contribution, retaining the explicitly bounded numerical contraction residual. Whole-head transfer used a same-field correct donor and can carry answer information; it was a diagnostic comparator, not a test of an answer-independent entity variable or a mathematical upper bound on repair performance.


Each change was computed once from the first-step capture and applied at the original query position on every full-sequence recomputation during generation. It did not move to the latest generated token or change with the generated answer. Each condition covered all 18 recipient facts, producing 162 complete answers. The diagnosis required at least four repairs across at least three worlds and at most one protected-answer failure for a repair branch to qualify. Whole-head nomination additionally required at least two more repairs than its matched random control. The random change was not norm-matched to the other factor changes and cannot establish their direction specificity.


\subsection*{Score diagnostics and statistical analysis}

The development sample size was fixed by the counterbalanced generator, not by a power calculation for a population-wide null effect. No behavioral responses were excluded for poor formatting, truncation, or low accuracy, and no outlier removal was applied. Causal cohorts were selected by the stated donor and wrong-object eligibility rule. Investigators were not blinded to the development labels or prior stage outcomes. Unsupported numerical inputs would stop the corresponding finite study rather than be replaced with more favorable cases; no unsupported conditions occurred in the reported validation or diagnosis.


We report counts with denominators, world-level variation, and the frozen 95\% descriptive percentile intervals for factual accuracy. The original routine used 2,000 resamples of whole worlds with random seed 20261114, retaining every repeat within each sampled world. These development intervals describe variation within the evaluated generator; they are not model-population confidence guarantees. We do not treat thousands of repeated questions as independent trials, report question-level binomial intervals, or interpret the six-world diagnosis as a population-wide absence result. No null-hypothesis significance tests were used for the primary claims.


Saved first-step score vectors supported an additional descriptive analysis without new model calls. For each original error, we measured the logit of the single canonical correct token minus that of the fixed first token of the original wrong response. We also calculated canonical-token probabilities over the full vocabulary and exact changes in the saved prediction vectors. This canonical-token margin is distinct from the accepted-variant score used during component discovery and from complete-answer accuracy. Continuous gains did not relax the failed practical criteria. Results were retained per question and world.


Additional wrong-object and eligible-pair intervals in S1 Data were calculated retrospectively using 2,000 shared whole-world resamples with seed 20261115. They are distinguished from the frozen factual-accuracy intervals and were not used for component selection or practical criteria.


Additional development-only checks realigned the name exchanges by queried word versus binding-line location and fitted simple predictive retrieval mixtures. The mixtures were fit on 24 worlds and evaluated on eight using four fixed whole-world folds. Expanded controls added constant object-identifier preferences and middle record positions to the original eight cues. These analyses were exploratory. Conditional wrong-object prediction removed the true object using labels and therefore was not a prospective predictor of whether a model would err. S1 Appendix specifies the comparisons and their interpretation limits.


\subsection*{Data and code availability}

Code and study data are available in the public repository at \url{https://github.com/Biodyn-AI/names-mechinterp} (version v0.1.0; release tag v0.1.0). Code is licensed under MIT; study data, documentation and paper are licensed under CC BY 4.0. S1 Data contains processed prompts, outcomes, figure coordinates and provenance records. Selected attention and component captures and historical audit receipts are included. The complete original raw native-logit safetensor archive is not included, so not every historical full-vocabulary audit can be repeated from this release. A Git tag is a version identifier, not a permanent archive; no archive DOI is claimed.


\subsection*{Reproducibility and use of artificial intelligence tools}

Every generation step retained its full-vocabulary native score vector. Separate CPU-only audits reconstructed greedy token choices, decoding, exact self trajectories, fixed query positions, component calculations, artifact hashes, and physical forward counts. An independently implemented grader reconstructed answers from the literal printed registry rather than reusing the experimental scoring routine. All reported numerical and whole-answer audits agreed. Internal plans and source closures were frozen before their respective execution stages; they were not public preregistrations or a journal Registered Report. The diagnosis and later score analyses were explicitly subsequent development analyses.


OpenAI Codex assisted experiment planning, implementation, execution orchestration, analysis, literature checks, figure-program development, and manuscript drafting. Output validity was assessed using deterministic tests, source and data hashes, complete saved-output replays, and separately implemented literal-answer checks. These automated checks do not substitute for a human author's assessment of the scientific claims. [Authors must confirm the exact tools and model versions and describe their own review and approval before submission.] No artificial intelligence system is listed as an author. The study used artificial registries and publicly released model weights; it involved no human participants, animals, or personal participant data.


\section*{Results}

\subsection*{Wrong object retrieval in both models}

Qwen answered 1,657 of 2,304 factual questions correctly, or 71.92\%, with a 95\% world-bootstrap interval of 64.71\% to 78.47\%. All 647 factual errors were completed answers giving another listed object's value in the requested field. Phi answered 1,980 of 2,304 correctly, or 85.94\%; its corresponding interval was 82.77\% to 88.93\%. Its errors included 260 wrong-object values and 64 other factual responses. Wrong-object errors occurred in 30 of 32 worlds for Qwen and 26 of 32 for Phi. Table 1 summarizes the outcome classes.


\begingroup\small\setstretch{1.15}
\begin{longtable}{|>{\raggedright\arraybackslash}p{\dimexpr0.4800\linewidth-2\tabcolsep-2\arrayrulewidth\relax}|>{\raggedright\arraybackslash}p{\dimexpr0.2600\linewidth-2\tabcolsep-2\arrayrulewidth\relax}|>{\raggedright\arraybackslash}p{\dimexpr0.2600\linewidth-2\tabcolsep-2\arrayrulewidth\relax}|}
\multicolumn{3}{@{}p{\linewidth}@{}}{\textbf{Table 1. Behavioral outcomes on the same 32 development worlds}}\\[5pt]
\hline
\textbf{Outcome} & \textbf{Qwen3 4B} & \textbf{Phi 3.5 mini}\\\hline
\endfirsthead
\hline
\textbf{Outcome} & \textbf{Qwen3 4B} & \textbf{Phi 3.5 mini}\\\hline
\endhead
All questions correct & 2297/3072 & 2257/3072\\\hline
Factual questions correct & 1657/2304 (71.92\%) & 1980/2304 (85.94\%)\\\hline
Fact accuracy world interval & 64.71\% to 78.47\% & 82.77\% to 88.93\%\\\hline
Wrong object fact answers & 647/2304 & 260/2304\\\hline
Other factual errors & 0/2304 & 64/2304\\\hline
Worlds with wrong object facts & 30/32 & 26/32\\\hline
Eligible equivalent name pairs & 117/768 & 69/768\\\hline
Worlds with eligible pairs & 23/32 & 14/32\\\hline
Strict identifier answers correct & 384/512 & 21/512\\\hline
Status answers correct & 256/256 & 256/256\\\hline
\end{longtable}
An eligible pair has one correct answer and one completed answer selecting another object's listed value. Other factual errors include nonlisted, malformed or truncated responses. Identifier scores enforce the requested bare answer format; Phi's low strict score must not be interpreted as absent identifier knowledge. Bootstrap intervals resample whole worlds.\par
\endgroup

Performance varied substantially by world. Qwen's factual accuracy ranged from 13/72 to 72/72, while Phi's ranged from 44/72 to 72/72. The two equivalent names disagreed in the eligible wrong-object sense in 117 of 768 Qwen pairs across 23 worlds and 69 of 768 Phi pairs across 14 worlds. Both models passed the specified feasibility criterion. Fig 1 displays world-level factual accuracy and eligible disagreement rates, making the repeated-measures structure visible rather than treating all questions as separate samples.


\begin{figure}[!h]\caption{Factual accuracy and equivalent name disagreement by world. Each point represents one of the same 32 worlds evaluated by both models. Panel A shows correct factual answers among 72 questions per world. Panel B shows eligible pairs among 24 equivalent-name pairs per world, requiring one correct answer and one completed wrong-object answer. Gray lines connect the same world between models. Horizontal offsets separate overlapping points and carry no additional variable. These are development observations, not independent model samples or held-out confirmation. Plotting code was developed with OpenAI Codex assistance; coordinates are calculated directly from the recorded outcomes, not generated illustrations.}\label{fig1}\end{figure}

Identifier responses required particular care. Qwen returned the correct bare identifier in 384/512 cases. Phi's strict score was 21/512, but an explicitly post hoc formatting inspection found 490 responses of the form Object followed by a valid identifier, of which 476 named the correct object. Together with the 21 correct bare identifiers, 497/512 responses were identifier-consistent under that secondary pattern check. The primary strict score was not changed. This distinction prevents a format failure from being presented as a failure to identify the referent.


\subsection*{Name and context controls did not yield one simple error rule}

Realigning the two-name exchange by word identity versus binding-line location did not reveal a single exact invariance pattern. Without highlighting, Qwen produced the same normalized response in all four cells of 70/96 world-by-field tables; 61 were consistently correct and nine consistently wrong. Four tables were name-aligned, seven line-aligned, and 15 matched neither exact signature. Phi had 81/96 all-same tables, two name-aligned, four line-aligned, and nine mixed tables. Some Phi differences concerned formatting rather than object choice. These categories describe output invariance, not uniquely identified internal retrieval routes.


Adding object-identifier and record-position preferences improved the global mixture's held-out predictive loss from 0.860 to 0.816 per factual response for Qwen and from 0.543 to 0.532 for Phi. The global mixtures nevertheless retained the correct binding as their most likely answer on the factual errors; the fit did not amount to a complete predictor of error occurrence. Richer field-conditioned controls and lexical response-word priors are reported with their full results in the accompanying data. Neither this limited baseline family nor the name-exchange signatures rules out the retrieval mechanisms described in prior work \cite{gurarieh2026mixing}.


\subsection*{Discovery effects did not establish a validated repair}

The early alias-boundary screen produced weak effects. Across its bounded head tests, first-token answer changes occurred in only one near-boundary discovery world. This did not establish that reference information was absent before the field query; it only failed to provide the desired repair at the screened locations.


At answer time, transferring a correct-name color profile at the ultimately selected head corrected the first predicted token in five of eight discovery cases. Its mean accepted-variant correct-versus-wrong score gain was 5.555, and its advantage over a cross-world profile control was 2.041. The head and operation were then frozen. On the six initially unused eligible development worlds, only one of eight complete-answer errors was repaired, in one world, while all ten initially correct recipient facts remained correct. The frozen average profile repaired none. Reverse and cyclic controls reached their specified wrong-object targets in only 1/18 questions each. Table 2 separates the stages and their outcome definitions.


\begingroup\small\setstretch{1.15}
\begin{longtable}{|>{\raggedright\arraybackslash}p{\dimexpr0.2000\linewidth-2\tabcolsep-2\arrayrulewidth\relax}|>{\raggedright\arraybackslash}p{\dimexpr0.2500\linewidth-2\tabcolsep-2\arrayrulewidth\relax}|>{\raggedright\arraybackslash}p{\dimexpr0.2800\linewidth-2\tabcolsep-2\arrayrulewidth\relax}|>{\raggedright\arraybackslash}p{\dimexpr0.2700\linewidth-2\tabcolsep-2\arrayrulewidth\relax}|}
\multicolumn{4}{@{}p{\linewidth}@{}}{\textbf{Table 2. Discovery and subsequent development tests}}\\[5pt]
\hline
\textbf{Stage} & \textbf{Worlds and question set} & \textbf{Observed outcome} & \textbf{Interpretation}\\\hline
\endfirsthead
\hline
\textbf{Stage} & \textbf{Worlds and question set} & \textbf{Observed outcome} & \textbf{Interpretation}\\\hline
\endhead
Answer time discovery & 8 worlds and 8 selected errors & 5/8 first tokens correct after transfer & Exploratory component selection\\\hline
Fixed repair validation & 6 other development worlds and 18 facts & 1/8 errors repaired and 10/10 correct retained & Repair criterion failed\\\hline
Reverse and cyclic controls & Same 6 worlds and 18 facts per control & 1/18 specified targets in each control & Specific steering criteria failed\\\hline
Fixed component diagnosis & Same reopened 6 worlds and 9 conditions & 162 complete answers and 360 forwards & All seven repair branches failed\\\hline
\end{longtable}
Discovery first-token success is not equivalent to completed-answer repair. The six-world diagnosis reused the failed validation cases and therefore adds a controlled explanation test, not an independent replication. The 64 reserved confirmation worlds were not evaluated.\par
\endgroup

\subsection*{Requested detail token attention remained concentrated}

Within the correct record, attention on the requested value's tokens averaged 94.98\% for the eight discovery errors and 96.12\% for the eight validation errors. Attention allocated to other field-value tokens was small. Total recipient attention to record tokens also remained high, averaging approximately 94.9\% in discovery and 95.8\% in validation. We observed no collapse in these attention-location measures. They do not establish which information the value vectors carried or rule out a wrong-field computation.


\subsection*{Complete component transfer also failed while prediction scores changed}

All 162 diagnostic answers completed within two tokens, including the final end-of-sequence token. The new same-field record-profile, within-record-pattern, combined-pattern, record-contribution, nonrecord-contribution, and whole-head conditions each repaired 0/8 errors. The original color-profile condition reproduced its earlier 1/8 repair exactly. Every condition preserved all ten initially correct answers. The random control repaired none. Table 3 reports all nine conditions without selecting a favorable subset.


\begingroup\small\setstretch{1.15}
\begin{longtable}{|>{\raggedright\arraybackslash}p{\dimexpr0.4100\linewidth-2\tabcolsep-2\arrayrulewidth\relax}|>{\raggedright\arraybackslash}p{\dimexpr0.1600\linewidth-2\tabcolsep-2\arrayrulewidth\relax}|>{\raggedright\arraybackslash}p{\dimexpr0.2000\linewidth-2\tabcolsep-2\arrayrulewidth\relax}|>{\raggedright\arraybackslash}p{\dimexpr0.2300\linewidth-2\tabcolsep-2\arrayrulewidth\relax}|}
\multicolumn{4}{@{}p{\linewidth}@{}}{\textbf{Table 3. Complete answer repairs and continuous score changes in the fixed diagnosis}}\\[5pt]
\hline
\textbf{Condition} & \textbf{Errors repaired} & \textbf{Correct facts retained} & \textbf{Mean token margin gain}\\\hline
\endfirsthead
\hline
\textbf{Condition} & \textbf{Errors repaired} & \textbf{Correct facts retained} & \textbf{Mean token margin gain}\\\hline
\endhead
Unchanged self control & 0/8 & 10/10 & 0.000\\\hline
Original color record profile & 1/8 & 10/10 & 5.996\\\hline
Same field record profile & 0/8 & 10/10 & 4.047\\\hline
Within record pattern only & 0/8 & 10/10 & 0.128\\\hline
Record profile and pattern & 0/8 & 10/10 & 4.194\\\hline
Whole record contribution & 0/8 & 10/10 & 4.207\\\hline
Nonrecord contribution & 0/8 & 10/10 & 0.003\\\hline
Whole donor head & 0/8 & 10/10 & 4.214\\\hline
Whole head matched random & 0/8 & 10/10 & 0.121\\\hline
\end{longtable}
Each row covers 18 factual questions in the same six development worlds. Margin gain is the mean change in canonical correct-token logit minus fixed original-wrong-token logit on the eight initial errors. It is not the discovery accepted-variant score or complete-answer accuracy. The random direction matches only the whole-head change's norm. Every nonself condition changed all 18 saved first-step score vectors; only the original color profile changed an answer.\par
\endgroup

Failure to repair did not mean that the intervention was numerically inactive. Whole-head transfer changed all 18 prediction vectors and improved the mean canonical correct-minus-original-wrong margin by 4.214 on the eight initial errors. The baseline mean margin was -10.904, so the average remained unfavorable. The whole-record contribution produced a similar gain of 4.207, whereas the nonrecord contribution gave a mean gain of 0.003. All new conditions left first-token choices unchanged. Fig 2 shows the individual error-level changes alongside complete-answer repairs. The two measures answer different questions: an intervention can causally favor the correct token without making the complete answer correct.


\begin{figure}[!h]\caption{Prediction score changes without reliable complete answer repair. Panel A shows canonical correct-minus-fixed-original-wrong logit changes for the eight initial factual errors in six worlds, with colors identifying worlds and a black marker giving the mean. Two errors from a world remain clustered observations. Panel B shows complete-answer repairs among the same eight errors. All conditions retained ten initially correct facts. Whole-head transfer and random directions are diagnostic comparisons, not performance bounds; only that random condition is norm-matched to the whole-head change. The reopened cases are development data. Plotting code was developed with OpenAI Codex assistance and draws directly from saved numerical outcomes.}\label{fig2}\end{figure}

All numerical replay checks and the independent literal-answer audit agreed. The 36 natural captures reproduced the original native prediction vectors exactly. All 18 unchanged-control answers and all 18 original color-transfer answers reproduced their earlier complete native trajectories. The maximum observed drift at the original query head was zero. CPU recomputation matched the saved component changes, and rescoring audited all 360 saved native endpoints without rerunning model inference. These checks support the validity of the negative outcome; they do not expand its scope beyond the fixed component, selected cases, and intervention families.


\section*{Discussion}

The main behavioral result is that equivalent names can produce different factual answers even when their common referent and its properties are explicitly defined in context. In both tested model families, the problem included completed answers drawn from another object's requested field, rather than only malformed responses or failure to recognize the answer vocabulary. The counterbalanced registries therefore provide a useful setting for separating reference selection from the ability to return a plausible field value. They do not establish that the models lack object representations or that every incorrect answer is caused by a name prior.


The causal results place a narrower constraint on one proposed explanation. The selected head had an informative attention pattern, and its modification exerted a measurable, directionally favorable influence on output scores. Nevertheless, the transferred profile repaired only one error in the six-world development test, and none of the new fixed-component variants repaired an error. Replacing more of the donor computation did not guarantee better performance: whole-head transfer repaired none, while the restricted color-profile operation repaired one. Component contributions can interact or cancel, so the whole-head condition should not be interpreted as a mathematical upper bound.


This finding fits the distinction between causal influence and a sufficient causal abstraction \cite{zhangnanda2024patching,geiger2025causal}. A component may contribute information relevant to the answer while other signals keep the wrong answer preferred. Our score changes demonstrate more than a uniform boost to all candidate answers: the correct token improved relative to the original wrong token. We therefore do not identify a generic vocabulary-boost mechanism. Nor does a positive margin change show that the selected head implements a transferable object-selection variable. The unsuccessful repair leaves several possibilities open, including cooperation among components, context-dependent downstream use, and more distributed reference computations.


The relationship to earlier binding work is consequently complementary rather than contradictory. Library-alias identification and attribute errors have already been dissociated \cite{yang2026import}. Lexical and positional retrieval routes can coexist \cite{gurarieh2026mixing}, and both retrieval-conditioned circuits and structured relational representations have been demonstrated elsewhere \cite{ohdemberg2026retrieval,dai2026cell}. Our results add a transparent development sequence in an arbitrary contextual registry: a promising intervention was fixed, failed on other cases, and remained unsuccessful under a targeted decomposition. They do not falsify those other mechanisms or make attention categorically uninformative \cite{jainwallace2019attention,wiegreffepinter2019attention}.


\subsection*{Limitations}

The experiment covers two small instruction-tuned models, eight alias words, two wording formats, and one synthetic task family. Words recur across worlds, and the design balances specified margins rather than every higher-order combination. The numerical inference configuration may differ from common reduced-precision deployments. Behavioral replication across model families does not constitute causal replication: only Qwen received internal interventions. Identifier formatting further limits broad identity-versus-property claims, especially for Phi.


The mechanistic cohorts were selected because they contained a strong correct donor and a completed wrong-object recipient. Eight errors across six worlds are a conditional development sample, not an estimate of universal repair failure. The diagnosis was designed after the initial validation outcome and reopened the same worlds. Reserved confirmation material was not evaluated. No family-wide absence claim, population-level zero-effect claim, or exhaustive search over components is supported. Preservation refers to the ten tested recipient facts; leaving other queries unmodified is not evidence of broad preservation across unrelated objects or tasks.


Same-field donor transfers can carry information about the requested answer, and the one random control is matched to only the whole-head norm. These comparisons diagnose parts of a fixed computation but do not establish direction-specific semantics for every factor. Canonical-token margins omit other accepted surface forms and cannot be substituted for complete-answer scoring. Internal execution freezes improve traceability but are not public preregistration. Finally, the current study supports a limited empirical contribution whose distinction from the closely related literature must be evaluated on its particular design and evidence, not on a claim of fundamental novelty.


\subsection*{Implications for further work}

A stronger entity-selection claim would require an operator chosen without consulting its confirmation outcomes, evaluated in fresh contexts with reassigned property values. A useful transfer should select the recipient object's own facts when those facts differ from the donor's, produce specified reverse effects, preserve relevant other answers, and replicate in another model family. Those are requirements for a future study, not unreported accomplishments of this one. The present negative result supplies a reason to stop treating the tested single head as an established portable repair interface.


\section*{Conclusions}

Equivalent names produced inconsistent factual retrieval in two models over the same 32 counterbalanced development worlds. In one model, an initially attractive reading-profile intervention failed its fixed development criteria, and a subsequent nine-condition diagnosis found no new complete-answer repair despite operative and often favorable score changes. The evidence supports reference inconsistency and the insufficiency of the tested repair hypothesis. The actual mechanism responsible for the errors remains unresolved.


\section*{Competing interests}

The author declares no competing interests.


\section*{Acknowledgments}

[Authors to confirm acknowledgments and obtain permission from any named individuals.]


\nolinenumbers\bibliography{references}

\section*{Supporting information}

S1 Appendix. Detailed methods and audit records. Extended task definitions, intervention construction, practical criteria, case-level outcomes, runtime identifiers, and reproducibility limits.


S1 Data. Processed outcomes and figure source data. Machine-readable prompts and behavioral answers for the 32 evaluated worlds, per-world summaries, fixed-component outcomes and score changes, source hashes and figure coordinates. This processed supplement is distinct from the release's selected native artifacts; the release-scope statement specifies which historical audits can be repeated.


\end{document}
